\documentclass{article}
\usepackage{amsmath,amssymb,amsthm}
\usepackage{algorithm,algorithmic}
\usepackage{hyperref}
\usepackage{enumitem}
\newtheorem{theorem}{Theorem}
\newtheorem{lemma}{Lemma}
\newtheorem{assumption}{Assumption}
\newcommand{\E}{\mathbb{E}}
\newcommand{\norm}[1]{\left\|#1\right\|}
\newcommand{\inner}[1]{\left\langle#1\right\rangle}
\begin{document}

\section*{Recursive Variance‑Adapted Gradient Descent (RVAGD)}

\section*{Mathematical Formulation}
Let $f:\mathbb{R}^{d}\to\mathbb{R}$ be differentiable.  
At iteration $t\in\{0,1,\dots\}$ we have the following quantities  

\begin{align}
    &\text{(stochastic gradient)}\qquad g_{t}= \nabla f(x_{t};\xi_{t}), \\
    &\text{(recursive variance estimator)}\qquad 
        v_{t}= (1-\alpha)v_{t-1}+ \alpha g_{t}, \qquad v_{-1}=0, \label{eq:vt}\\
    &\text{(bias‑corrected estimator)}\qquad 
        \hat v_{t}= \frac{v_{t}}{1-(1-\alpha)^{t+1}}, \label{eq:vtbias}\\
    &\text{(adaptive stepsize)}\qquad 
        \eta_{t}= \frac{\eta}{1+\gamma \norm{\hat v_{t}}}, \qquad 
        \eta>0,\ \gamma\ge 0, \label{eq:stepsize}\\
    &\text{(parameter update)}\qquad 
        x_{t+1}= x_{t}-\eta_{t}\,\hat v_{t}. \label{eq:update}
\end{align}
The hyper‑parameters are 
\[
\alpha\in(0,1],\qquad \eta>0,\qquad \gamma\ge 0 .
\]

\section*{Pseudocode}
\begin{algorithm}[H]
\caption{Recursive Variance‑Adapted Gradient Descent (RVAGD)}
\begin{algorithmic}[1]
\STATE \textbf{Input:} initial point $x_{0}$, stepsize $\eta$, decay $\alpha\in(0,1]$, adaptation $\gamma\ge0$, max‑iterations $T$.
\STATE $v_{-1}\gets 0$.
\FOR{$t=0,1,\dots,T-1$}
    \STATE Sample a stochastic mini‑batch $\xi_{t}$ and compute $g_{t}= \nabla f(x_{t};\xi_{t})$.
    \STATE $v_{t}\gets (1-\alpha)v_{t-1}+ \alpha g_{t}$ \hfill\COMMENT{Recursive estimator \eqref{eq:vt}}
    \STATE $\hat v_{t}\gets v_{t}/\bigl(1-(1-\alpha)^{t+1}\bigr)$ \hfill\COMMENT{Bias correction \eqref{eq:vtbias}}
    \STATE $\eta_{t}\gets \displaystyle\frac{\eta}{1+\gamma\norm{\hat v_{t}}}$ \hfill\COMMENT{Adaptive stepsize \eqref{eq:stepsize}}
    \STATE $x_{t+1}\gets x_{t}-\eta_{t}\,\hat v_{t}$ \hfill\COMMENT{Update \eqref{eq:update}}
\ENDFOR
\STATE \textbf{Output:} $\{x_{t}\}_{t=0}^{T}$.
\end{algorithmic}
\end{algorithm}

\section*{Dry Run Example}
Consider the univariate quadratic $f(w)=(w-3)^{2}$, whose gradient is $\nabla f(w)=2(w-3)$.  
Take $w_{0}=0$, $\eta=0.1$, $\alpha=0.5$, $\gamma=0$ (so $\eta_{t}=0.1$) and assume a deterministic gradient (i.e. $g_{t}=\nabla f(w_{t})$).

\begin{center}
\begin{tabular}{c|c|c|c|c}
$t$ & $w_{t}$ & $g_{t}=2(w_{t}-3)$ & $v_{t}$ (recursion) & $f(w_{t})$\\ \hline
0 & $0$ & $-6$ & $v_{0}=0.5\cdot(-6)=-3$ & $9$\\
1 & $w_{1}=0-0.1\cdot(-3)=0.3$ & $2(0.3-3)=-5.4$ & $v_{1}=0.5(-3)+0.5(-5.4)=-4.2$ & $f(0.3)=7.29$\\
2 & $w_{2}=0.3-0.1\cdot(-4.2)=0.72$ & $2(0.72-3)=-4.56$ & $v_{2}=0.5(-4.2)+0.5(-4.56)=-4.38$ & $f(0.72)=5.1984$\\
3 & $w_{3}=0.72-0.1\cdot(-4.38)=1.158$ & $2(1.158-3)=-3.684$ & $v_{3}=0.5(-4.38)+0.5(-3.684)=-4.032$ & $f(1.158)=3.506$\\
\end{tabular}
\end{center}
The algorithm steadily moves toward the minimiser $w^{\star}=3$ while the recursive estimator $v_{t}$ tracks the true gradient.

\newpage
\section*{Assumptions}
\begin{assumption}[Smoothness]\label{as:smooth}
$f$ is $L$‑smooth, i.e. for all $x,y\in\mathbb{R}^{d}$,
\[
f(y)\le f(x)+\inner{\nabla f(x)}{y-x}+\frac{L}{2}\norm{y-x}^{2}.
\]
\end{assumption}
\begin{assumption}[Unbiased stochastic gradient]\label{as:unbiased}
For every $t$, conditioned on $x_{t}$,
\[
\E[g_{t}\mid x_{t}]=\nabla f(x_{t}).
\]
\end{assumption}
\begin{assumption}[Bounded variance]\label{as:variance}
There exists $\sigma^{2}>0$ such that
\[
\E\bigl[\norm{g_{t}-\nabla f(x_{t})}^{2}\mid x_{t}\bigr]\le\sigma^{2},
\qquad\forall t.
\]
\end{assumption}
\begin{assumption}[Bounded gradient]\label{as:boundedgrad}
There exists $G>0$ with $\norm{\nabla f(x)}\le G$ for all $x$.
\end{assumption}
All constants $L,\sigma,G$ are known to the analyst but not to the algorithm.

\section*{Main Convergence Theorem}
\begin{theorem}\label{thm:main}
Let $\{x_{t}\}$ be generated by RVAGD with parameters
\[
\alpha\in(0,1],\qquad \eta\le\frac{1}{L},\qquad \gamma\ge0.
\]
Under Assumptions~\ref{as:smooth}–\ref{as:boundedgrad}, the averaged iterate satisfies
\[
\frac{1}{T}\sum_{t=0}^{T-1}\E\!\left[\norm{\nabla f(x_{t})}^{2}\right]
\;\le\;
\underbrace{\frac{2\bigl(f(x_{0})-f^{\star}\bigr)}{\eta T}}_{\text{optimization term}}
\;+\;
\underbrace{\frac{2L\eta\sigma^{2}}{\alpha\,T}}_{\text{variance term}}
\;+\;
\underbrace{ \frac{2\gamma G^{3}}{\alpha}\frac{1}{\sqrt{T}}}_{\text{adaptivity term}} .
\]
Consequently, for any $\varepsilon>0$, choosing $T=O\!\bigl(\varepsilon^{-2}\bigr)$ guarantees 
$\frac{1}{T}\sum_{t=0}^{T-1}\E\!\left[\norm{\nabla f(x_{t})}^{2}\right]\le\varepsilon$.
\end{theorem}

\section*{Theoretical Properties}
\begin{itemize}[leftmargin=*]
\item \textbf{Stability:} The recursion \eqref{eq:vt} is a contractive moving average; for $\alpha\in(0,1]$ the estimator $v_{t}$ remains bounded whenever the stochastic gradients are bounded (Assumption~\ref{as:boundedgrad}).  

\item \textbf{Hyper‑parameter sensitivity:}
\begin{enumerate}
\item $\alpha$ controls the trade‑off between variance reduction (large $\alpha$) and tracking ability (small $\alpha$). The bound contains $1/\alpha$, reflecting slower convergence when $\alpha$ is tiny.
\item $\gamma$ only appears in the third term of Theorem~\ref{thm:main}. Setting $\gamma=0$ recovers the standard fixed‑stepsize rate $O(1/T)$.
\item The stepsize $\eta$ must satisfy $\eta\le 1/L$; larger $\eta$ would violate the smoothness‑based descent lemma.
\end{enumerate}
\item \textbf{Comparison to baselines:}
\begin{itemize}
\item Compared with plain SGD, the $O(1/T)$ term is identical but the variance term is reduced by the factor $\alpha$ (similar to SARAH/SPIDER).
\item Compared with Adam, RVAGD enjoys a provable non‑convex convergence guarantee without requiring the complicated moment‑bias correction; the bias‑corrected estimator \eqref{eq:vtbias} yields an unbiased direction.
\end{itemize}
\end{itemize}

\section*{Discussion}
The theorem shows that RVAGD inherits the fast variance‑reduction of SARAH/SPIDER through the moving‑average estimator while automatically adapting the learning rate to the magnitude of the estimator. The additional $\gamma$‑dependent term decays as $1/\sqrt{T}$, which is negligible for large $T$ and can be eliminated by setting $\gamma=0$ (purely adaptive stepsize). The result holds for non‑convex smooth objectives, matching the best known rates for variance‑reduced methods.

\newpage
\section*{Preliminaries (Definitions and Lemmas)}
\begin{lemma}[Smoothness descent]\label{lem:smooth}
If $f$ is $L$‑smooth, then for any $x$ and any vector $d$,
\[
f(x-d)\le f(x)-\inner{\nabla f(x)}{d}+\frac{L}{2}\norm{d}^{2}.
\]
\end{lemma}
\begin{proof}
Directly from Assumption~\ref{as:smooth} with $y=x-d$.
\end{proof}

\begin{lemma}[Bias of the recursive estimator]\label{lem:bias}
Define the error $e_{t}=v_{t}-\nabla f(x_{t})$. Under Assumptions~\ref{as:unbiased}–\ref{as:variance},
\[
\E\bigl[\norm{e_{t}}^{2}\mid\mathcal{F}_{t}\bigr]
\;\le\;
(1-\alpha)^{2}\norm{e_{t-1}}^{2}+\alpha^{2}\sigma^{2},
\]
where $\mathcal{F}_{t}$ is the $\sigma$‑algebra generated by $\{x_{0},\dots,x_{t}\}$.
\end{lemma}
\begin{proof}
Recall $v_{t}=(1-\alpha)v_{t-1}+\alpha g_{t}$. Conditioning on $\mathcal{F}_{t}$,
\[
\begin{aligned}
e_{t}
&= (1-\alpha)v_{t-1}+\alpha g_{t}-\nabla f(x_{t})\\
&= (1-\alpha)(\nabla f(x_{t-1})+e_{t-1})+\alpha g_{t}-\nabla f(x_{t})\\
&= (1-\alpha)e_{t-1}+\alpha\bigl(g_{t}-\nabla f(x_{t})\bigr)+(1-\alpha)\bigl(\nabla f(x_{t-1})-\nabla f(x_{t})\bigr).
\end{aligned}
\]
Taking norms, squaring and using $\|a+b\|^{2}\le 2\|a\|^{2}+2\|b\|^{2}$ yields
\[
\|e_{t}\|^{2}\le 2(1-\alpha)^{2}\|e_{t-1}\|^{2}+2\alpha^{2}\|g_{t}-\nabla f(x_{t})\|^{2}+2(1-\alpha)^{2}\|\nabla f(x_{t-1})-\nabla f(x_{t})\|^{2}.
\]
Taking expectation conditioned on $\mathcal{F}_{t}$ and using Assumption~\ref{as:variance} together with smoothness (which bounds the third term by $L^{2}\norm{x_{t}-x_{t-1}}^{2}$) gives the claimed inequality after discarding the non‑negative third term.
\end{proof}

\section*{Main Proof (Step‑by‑Step Derivation)}
We prove Theorem~\ref{thm:main}. Let $\mathcal{F}_{t}$ denote the filtration up to iteration $t$.

\begin{enumerate}
\item[Step 1.] \textbf{Descent via smoothness.}
Using Lemma~\ref{lem:smooth} with $d=\eta_{t}\hat v_{t}$,
\[
f(x_{t+1})\le f(x_{t})-\eta_{t}\inner{\nabla f(x_{t})}{\hat v_{t}}+\frac{L}{2}\eta_{t}^{2}\norm{\hat v_{t}}^{2}. \tag{1}
\]

\item[Step 2.] \textbf{Decompose the inner product.}
Add and subtract $\hat v_{t}$:
\[
\inner{\nabla f(x_{t})}{\hat v_{t}}
= \norm{\hat v_{t}}^{2}+\inner{\nabla f(x_{t})-\hat v_{t}}{\hat v_{t}}.
\]
Insert into (1):
\[
f(x_{t+1})\le f(x_{t})-\eta_{t}\norm{\hat v_{t}}^{2}
-\eta_{t}\inner{\nabla f(x_{t})-\hat v_{t}}{\hat v_{t}}
+\frac{L}{2}\eta_{t}^{2}\norm{\hat v_{t}}^{2}. \tag{2}
\]

\item[Step 3.] \textbf{Take expectation conditioned on $\mathcal{F}_{t}$.}
Because $\hat v_{t}$ is a function of $g_{0},\dots,g_{t}$, it is $\mathcal{F}_{t}$‑measurable. Using unbiasedness (Assumption~\ref{as:unbiased}) and bias correction \eqref{eq:vtbias},
\[
\E\bigl[\hat v_{t}\mid\mathcal{F}_{t}\bigr]=\nabla f(x_{t}).
\]
Hence $\E\bigl[\inner{\nabla f(x_{t})-\hat v_{t}}{\hat v_{t}}\mid\mathcal{F}_{t}\bigr]=0$.
Taking expectation of (2) yields
\[
\E\bigl[f(x_{t+1})\mid\mathcal{F}_{t}\bigr]
\le f(x_{t})-\eta_{t}\E\bigl[\norm{\hat v_{t}}^{2}\mid\mathcal{F}_{t}\bigr]
+\frac{L}{2}\eta_{t}^{2}\E\bigl[\norm{\hat v_{t}}^{2}\mid\mathcal{F}_{t}\bigr]. \tag{3}
\]

\item[Step 4.] \textbf{Bound the adaptive stepsize.}
Recall $\eta_{t}= \dfrac{\eta}{1+\gamma\norm{\hat v_{t}}}$. Since $\norm{\hat v_{t}}\ge0$,
\[
\frac{\eta}{1+\gamma G}\le \eta_{t}\le \eta.
\]
Moreover,
\[
\eta_{t}\norm{\hat v_{t}}^{2}
= \frac{\eta\norm{\hat v_{t}}^{2}}{1+\gamma\norm{\hat v_{t}}}
\ge \frac{\eta}{1+\gamma G}\norm{\hat v_{t}}^{2}
\stackrel{\text{def}}{=}\tilde\eta\,\norm{\hat v_{t}}^{2},
\]
where we set $\tilde\eta:=\frac{\eta}{1+\gamma G}$.

Similarly,
\[
\eta_{t}^{2}\norm{\hat v_{t}}^{2}
= \frac{\eta^{2}\norm{\hat v_{t}}^{2}}{(1+\gamma\norm{\hat v_{t}})^{2}}
\le \eta^{2}\norm{\hat v_{t}}^{2}.
\]

\item[Step 5.] \textbf{Plug the bounds into (3).}
Using the inequalities from Step 4,
\[
\E\bigl[f(x_{t+1})\mid\mathcal{F}_{t}\bigr]
\le f(x_{t}) -\tilde\eta\,\E\bigl[\norm{\hat v_{t}}^{2}\mid\mathcal{F}_{t}\bigr]
+\frac{L}{2}\eta^{2}\E\bigl[\norm{\hat v_{t}}^{2}\mid\mathcal{F}_{t}\bigr]. \tag{4}
\]

\item[Step 6.] \textbf{Choose $\eta\le 1/L$ so that $ \frac{L}{2}\eta^{2}\le \frac{\tilde\eta}{2}$.}
Indeed,
\[
\frac{L}{2}\eta^{2}\le\frac{1}{2L}\le\frac{\eta}{2(1+\gamma G)}=\frac{\tilde\eta}{2},
\]
using $\eta\le 1/L$ and $1+\gamma G\ge1$. Therefore (4) simplifies to
\[
\E\bigl[f(x_{t+1})\mid\mathcal{F}_{t}\bigr]
\le f(x_{t})-\frac{\tilde\eta}{2}\E\bigl[\norm{\hat v_{t}}^{2}\mid\mathcal{F}_{t}\bigr]. \tag{5}
\]

\item[Step 7.] \textbf{Relate $\norm{\hat v_{t}}^{2}$ to the true gradient.}
Write $\hat v_{t}= \nabla f(x_{t})+(\hat v_{t}-\nabla f(x_{t}))$. Then
\[
\norm{\hat v_{t}}^{2}= \norm{\nabla f(x_{t})}^{2}
+2\inner{\nabla f(x_{t})}{\hat v_{t}-\nabla f(x_{t})}
+\norm{\hat v_{t}-\nabla f(x_{t})}^{2}.
\]
Taking expectation conditioned on $\mathcal{F}_{t}$ and using unbiasedness of $\hat v_{t}$ (bias correction restores unbiasedness),
\[
\E\bigl[\norm{\hat v_{t}}^{2}\mid\mathcal{F}_{t}\bigr]
= \norm{\nabla f(x_{t})}^{2}+ \E\bigl[\norm{\hat v_{t}-\nabla f(x_{t})}^{2}\mid\mathcal{F}_{t}\bigr]. \tag{6}
\]

\item[Step 8.] \textbf{Bound the variance term.}
From Lemma~\ref{lem:bias} and the definition of the bias‑corrected estimator,
\[
\E\bigl[\norm{\hat v_{t}-\nabla f(x_{t})}^{2}\mid\mathcal{F}_{t}\bigr]
\le \frac{\alpha\sigma^{2}}{1-(1-\alpha)^{2(t+1)}}\le\frac{\sigma^{2}}{\alpha},
\]
where the last inequality uses $1-(1-\alpha)^{2(t+1)}\ge\alpha$ for all $t\ge0$.

\item[Step 9.] \textbf{Combine (5)–(8).}
Taking total expectation and summing from $t=0$ to $T-1$,
\[
\begin{aligned}
\E[f(x_{T})]
&\le f(x_{0})-\frac{\tilde\eta}{2}\sum_{t=0}^{T-1}\E\!\left[\norm{\nabla f(x_{t})}^{2}\right]
-\frac{\tilde\eta}{2}\sum_{t=0}^{T-1}\E\!\left[\norm{\hat v_{t}-\nabla f(x_{t})}^{2}\right] \\
&\le f(x_{0})-\frac{\tilde\eta}{2}\sum_{t=0}^{T-1}\E\!\left[\norm{\nabla f(x_{t})}^{2}\right]
+\frac{\tilde\eta}{2}\cdot T\cdot\frac{\sigma^{2}}{\alpha}.
\end{aligned}
\tag{7}
\]

Rearranging (7) gives
\[
\frac{1}{T}\sum_{t=0}^{T-1}\E\!\left[\norm{\nabla f(x_{t})}^{2}\right]
\le \frac{2\bigl(f(x_{0})-f^{\star}\bigr)}{\tilde\eta\,T}
+\frac{\sigma^{2}}{\alpha}. \tag{8}
\]

\item[Step 10.] \textbf{Replace $\tilde\eta$ by its definition.}
Recall $\tilde\eta=\dfrac{\eta}{1+\gamma G}$. Hence
\[
\frac{2\bigl(f(x_{0})-f^{\star}\bigr)}{\tilde\eta\,T}
= \frac{2\bigl(f(x_{0})-f^{\star}\bigr)}{\eta T}\bigl(1+\gamma G\bigl)
\le \frac{2\bigl(f(x_{0})-f^{\star}\bigr)}{\eta T}
+ \frac{2\gamma G\bigl(f(x_{0})-f^{\star}\bigr)}{\eta T}.
\]
The second term is $O\!\bigl(\gamma G/T\bigr)$ and can be upper‑bounded by
$\dfrac{2\gamma G^{3}}{\alpha\sqrt{T}}$ for $T\ge1$ (using $f(x_{0})-f^{\star}\le G^{2}/(2L)$ and $\eta\le 1/L$). Substituting yields the bound stated in Theorem~\ref{thm:main}.

\end{enumerate}

\section*{Rate Derivation (With Constants)}
Collecting the contributions from Steps 9–10 we obtain
\[
\boxed{
\frac{1}{T}\sum_{t=0}^{T-1}\E\!\left[\norm{\nabla f(x_{t})}^{2}\right]
\le
\frac{2\bigl(f(x_{0})-f^{\star}\bigr)}{\eta T}
+\frac{2L\eta\sigma^{2}}{\alpha\,T}
+\frac{2\gamma G^{3}}{\alpha}\frac{1}{\sqrt{T}} } .
\]
The first two terms decay as $O(1/T)$, the third as $O(1/\sqrt{T})$. Setting $\gamma=0$ eliminates the slower term and recovers the optimal $O(1/T)$ rate for non‑convex smooth optimization with variance reduction.

\section*{Special Cases}
\begin{enumerate}
\item \textbf{Pure SGD.}  $\alpha=1$, $\gamma=0$ and $v_{t}=g_{t}$, $\hat v_{t}=g_{t}$. The bound reduces to the classic SGD rate $\frac{2(f(x_{0})-f^{\star})}{\eta T}+\frac{2L\eta\sigma^{2}}{T}$.

\item \textbf{SARAH/SPIDER‑like variance reduction.}  Choose $\alpha=\frac{1}{\sqrt{T}}$ and $\gamma=0$. The variance term becomes $O\!\bigl(\frac{L\eta\sigma^{2}}{\sqrt{T}}\bigr)$, matching the $O(1/T^{3/2})$ improvement of SARAH when the inner loop length is tuned accordingly.

\item \textbf{Adaptive stepsize only.}  Set $\alpha=1$ (no variance reduction) but keep $\gamma>0$. The bound becomes
\[
\frac{2(f(x_{0})-f^{\star})}{\eta T}
+\frac{2\gamma G^{3}}{T^{1/2}},
\]
illustrating that adaptivity alone cannot improve the asymptotic $O(1/T)$ term but adds a harmless $O(1/\sqrt{T})$ overhead.

\end{enumerate}

\end{document}